\documentclass[11pt]{article}
\usepackage[margin=1.3in]{geometry}
\usepackage[utf8]{inputenc}
\usepackage[UKenglish]{babel}

\usepackage{amsmath}
\usepackage{amsthm}
\usepackage{amsfonts}
\usepackage{amssymb}
%\usepackage{nath}

\newtheorem{theorem}{Theorem}
\newtheorem{lemma}{Lemma}
\newtheorem{corollary}{Corollary}
\newtheorem{definition}{Definition}
\newtheorem{example}{Example}
\newtheorem{remark}{Remark}
\newtheorem{fact}{Fact}


\usepackage{algorithm}
\usepackage[noend]{algpseudocode}
% For using [H] float placement in algorithms
\usepackage{float}

% Thomas packages


\usepackage{todonotes}
\presetkeys{todonotes}{inline}{}

\usepackage{enumitem}
\usepackage{multicol}
\usepackage[hypertexnames=false]{hyperref}
\makeatletter
% Avoid duplicate hyperlink destinations when counters reset (e.g. algorithmic line numbers).
\providecommand*{\theHALG@line}{\arabic{algorithm}.\arabic{ALG@line}}
\makeatother
\usepackage{cleveref}
\usepackage{autonum}
\DeclareMathAlphabet{\mathpzc}{OT1}{pzc}{m}{it}
% Figures
\usepackage{graphicx}
\usepackage{tikz}
\usetikzlibrary{arrows.meta,positioning,calc}

\DeclareMathOperator{\poly}{poly}
\DeclareMathOperator{\sign}{sign}
\DeclareMathOperator{\asvec}{vec}
\DeclareMathOperator*{\argmax}{arg\,max}
\DeclareMathOperator*{\argmin}{arg\,min}
% Allow page break in align
\allowdisplaybreaks


\usepackage{mathtools}
\DeclarePairedDelimiterX{\infdivx}[2]{(}{)}{#1 \,\delimsize\|\, #2}
\DeclareMathOperator*{\Div}{D}
\newcommand{\D}{\Div\infdivx}
\DeclareMathOperator*{\smallDiv}{d}
\newcommand{\di}{\smallDiv\infdivx}
\DeclareMathOperator*{\h}{h}

% Jakob packages
\newcommand{\eps}{\varepsilon}
\newcommand{\R}{\mathbb{R}}
\newcommand{\Z}{\mathbb{Z}}
\newcommand{\N}{\mathbb{N}}
\newcommand{\F}{\mathbb{F}}
\newcommand{\C}{\mathbb{C}}
\newcommand{\Q}{\mathbb{Q}}

\newcommand{\abs}[1]{\left\lvert #1\right\rvert}
\newcommand{\set}[1]{\left\{ #1 \right\}}
\newcommand{\setbuilder}[2]{\left\{ #1 \; \middle\vert \; #2 \right\}}
\newcommand{\norm}[2]{\left\| #1 \right\|_{#2}}
\newcommand{\innerprod}[2]{\left\langle #1, \; #2 \right\rangle}

\newcommand{\Prp}[1]{\Pr\!\left[#1 \right]}
\newcommand{\Prpcond}[2]{\Pr\!\left[#1 \; \middle| \; #2 \right]}
\DeclareMathOperator*{\E}{E}
\newcommand{\Epcond}[2]{\E\!\left[#1 \; \middle| \; #2 \right]}

\newcommand{\br}[1]{\left( #1 \right)}


\newcommand{\smat}[1]{\left[\begin{smallmatrix}#1\end{smallmatrix}\right]}
\newcommand{\svec}[1]{[\begin{smallmatrix}#1\end{smallmatrix}]}

\newcommand{\floor}[1]{\left\lfloor {#1} \right\rfloor}

% Index sets (avoid ambiguity of [n] across notes)
% Inclusive: {0,1,...,n}
\newcommand{\idx}[1]{\{0,\dotsc,#1\}}
% Half-open: {0,1,...,n-1}
\newcommand{\rng}[1]{\{0,\dotsc,#1-1\}}

% Better math spacing in coefficient extractor
\newcommand{\coeff}[2]{[x^{#2}]\,#1}
% Clarify [n] notation inline where needed
% \newcommand{\range}[1]{\{0,\ldots,#1\}}
